% Generated by scripts/assemble.py; corrections belong in source/corrections.json.
\examhead{188}{陕西师范大学 2025 年数学分析试卷}{177}
% p178-b13; source PDF page 178
\par\Needspace{5\baselineskip}
\blockanchor{p178-b13}
\majorquestion{一}计算题. 每题 5 分, 共 30 分.\par
% p178-b14; source PDF page 178
\blockanchor{p178-b14}
\question{1}求极限  $\lim_{n\to\infty}\frac{1!+2!+\cdots+n!}{n!}$ .\par
% p178-b15; source PDF page 178
\blockanchor{p178-b15}
\question{2}求极限 $\lim_{x\to +\infty}x^2 [\ln \arctan (x + 1) - \ln \arctan x]$.\par
% p178-b16; source PDF page 178
\blockanchor{p178-b16}
\question{3}设 $f(x) = \int_{0}^{x} \frac{\sin t}{\pi - t} \, \mathrm{d}t$, 求积分 $\int_{0}^{\pi} f(x) \, \mathrm{d}x$.\par
% p178-b17; source PDF page 178
\blockanchor{p178-b17}
\question{4}计算 $\iint_{S} \frac{ax \, \mathrm{d}y \, \mathrm{d}z + (a + z)^2 \, \mathrm{d}x \, \mathrm{d}y}{\sqrt{x^2 + y^2 + z^2}} (a > 0)$, $S$ 取半球面 $z = -\sqrt{a^2 - x^2 - y^2}$ 的上侧.\par
% p178-b18; source PDF page 178
\blockanchor{p178-b18}
\question{5}计算 $I = \oint_{L} \frac{x \, \mathrm{d}y - y \, \mathrm{d}x}{3x^2 + 4y^2}$，其中 $L$ 是椭圆 $2x^2 + 3y^2 = 1$，方向沿逆时针方向.\par
% p178-b19; source PDF page 178
\blockanchor{p178-b19}
\question{6}求数项级数 $\sum_{n=1}^{\infty} \frac{(-1)^{n+1}}{n(n+1)}$ 的和.\par
% p178-b20; source PDF page 178
\par\Needspace{5\baselineskip}
\blockanchor{p178-b20}
\majorquestion{二}解答题. 每题 10 分, 共 30 分.\par
% p178-b21; source PDF page 178
\blockanchor{p178-b21}
\question{7}设  $f(x,y)=\sqrt[3]{x^{3}+y^{3}}$ .\par
% p178-b22; source PDF page 178
\blockanchor{p178-b22}
(1) 讨论 $f(x, y)$ 在原点处沿任何方向的导数是否均存在？\par
% p178-b23; source PDF page 178
\blockanchor{p178-b23}
(2) 讨论 $f_{x}(0,0), f_{y}(0,0)$ 是否存在？若存在，求出其值.\par
% p178-b24; source PDF page 178
\blockanchor{p178-b24}
(3) 讨论 $f(x, y)$ 在 $(0, 0)$ 处是否可微？\par
% p178-b25; source PDF page 178
\blockanchor{p178-b25}
\question{8}求方程 $2x^{2} + 2y^{2} + z^{2} + 8zy - z + 8 = 0$ 所确定的隐函数 $z = z(x, y)$ 的极值.\par
% p178-b26; source PDF page 178
\blockanchor{p178-b26}
\question{9}设 $u_{n} \neq 0$, 且 $\lim_{n \to \infty} \frac{n}{u_n} = 1$. 讨论级数 $\sum_{n=1}^{\infty} (-1)^{n+1} \left( \frac{1}{u_n} + \frac{1}{u_{n+1}} \right)$ 是否绝对收敛.\par
% p179-b00; source PDF page 179
\par\Needspace{5\baselineskip}
\blockanchor{p179-b00}
\majorquestion{三}证明题. 每题 15 分, 共 90 分.\par
% p179-b01; source PDF page 179
\blockanchor{p179-b01}
\question{10}设 $f(x)$ 在 $[a, +\infty)$ 上连续, 当 $x \to +\infty$ 时, $y = 2x$ 是 $f(x)$ 的渐近线. 证明: $f(x)$ 在 $[a, +\infty)$ 上一致连续.\par
% p179-b02; source PDF page 179
\blockanchor{p179-b02}
\question{11}设函数 $f(x)$ 在 $[0,1]$ 上二阶可导，且 $f''(x) \geq 0$. 证明：$\int_0^1 f(x^6) \mathrm{d}x \geq f\left(\frac{1}{7}\right)$.\par
% p179-b03; source PDF page 179
\blockanchor{p179-b03}
\question{12}设 $f(x)$ 在 $[0, +\infty)$ 上单调减, 且 $\lim_{x \to +\infty} f(x) = 0$. 证明: 若 $f'(x)$ 在 $[0, +\infty)$ 上连续, 则反常积分 $\int_0^{+\infty} f'(x) \sin^2 x \, \mathrm{d}x$ 收敛.\par
% p179-b04; source PDF page 179
\blockanchor{p179-b04}
\question{13}设 $f(x) = \sum_{n=1}^{\infty} \frac{\cos nx}{\sqrt{n^3 + n}}$.\par
% p179-b05; source PDF page 179
\blockanchor{p179-b05}
(1) 证明: $f(x)$ 在 $(- \infty, + \infty)$ 上连续.\par
% p179-b06; source PDF page 179
\blockanchor{p179-b06}
(2) 证明: $F\left(\frac{\pi}{2}\right) < \frac{\sqrt{2}}{2}$, 其中 $F(x) = \int_{0}^{x} f(t) \, \mathrm{d}t$.\par
% p179-b07; source PDF page 179
\blockanchor{p179-b07}
\question{14}证明 $I(y) = \int_{0}^{+\infty} e^{-px} \cos xy \, \mathrm{d}x$ 在 $(- \infty, +\infty)$ 上一致收敛, 其中 $p > 0$.\par
% p179-b08; source PDF page 179
\blockanchor{p179-b08}
\question{15}解答如下问题:\par
% p179-b09; source PDF page 179
\blockanchor{p179-b09}
(1) 计算积分 $A = \int_{0}^{1} \int_{0}^{1} \left| xy - \frac{1}{4} \right| \mathrm{d}x \mathrm{d}y$.\par
% p179-b10; source PDF page 179
\par\Needspace{6\baselineskip}
\blockanchor{p179-b10}
(2) 设 $z = f(x, y)$ 在区域 $D: [0, 1] \times [0, 1]$ 上连续，且满足条件\par
\nopagebreak[4]
% p179-b11; source PDF page 179
\blockanchor{p179-b11}
\displayformula{\iint_ {D} f (x, y)   \mathrm{d} x \mathrm{d} y = 0   \text {和}   \iint_ {D} x y f (x, y)   \mathrm{d} x \mathrm{d} y = 1.}
% p179-b12; source PDF page 179
\blockanchor{p179-b12}
证明: 存在  $(\xi,\eta)\in D$ ，使得  $|f(\xi,\eta)|\geq\frac{1}{A}$ ，其中 A 为 (1) 中的积分值.\par
